% Generated draft: requires mathematical and source review.
\documentclass[11pt]{article}
\usepackage[margin=25mm]{geometry}
\usepackage{fontspec}
\setmainfont{Latin Modern Roman}
\usepackage{amsmath,amssymb,mathtools,unicode-math}
\setmathfont{Latin Modern Math}
\usepackage{enumitem,needspace,xurl,hyperref}
\setlength{\parindent}{0pt}
\setlength{\parskip}{5pt}
\setlength{\emergencystretch}{3em}
\newcommand{\N}{\mathbb{N}}
\newcommand{\Z}{\mathbb{Z}}
\newcommand{\Q}{\mathbb{Q}}
\newcommand{\R}{\mathbb{R}}
\newcommand{\C}{\mathbb{C}}
\newcommand{\F}{\mathbb{F}}
\newcommand{\A}{\mathbb{A}}
\newcommand{\PP}{\mathbb{P}}
\newcommand{\E}{\mathbb{E}}
\newcommand{\eps}{\varepsilon}
\newcommand{\dd}{\,\mathrm{d}}
\begin{document}
\section*{UnsolvedMath ID: 700046}
\textbf{Source problem code:} catalogue entry
\subsection{Definitions and mathematical statement}
For a field \(K\) and a positive integer \(n\), let \(\operatorname{GL}_n(K)\) denote the group of invertible \(n\times n\) matrices over \(K\). A matrix group over \(K\) means a subgroup \(G\leq\operatorname{GL}_n(K)\).

A strict total order \(<\) on a group \(G\) is left invariant if
\[
\forall g,x,y\in G,\qquad x<y\Longrightarrow gx<gy.
\]
It is invariant on both sides if
\[
\forall a,b,x,y\in G,\qquad x<y\Longrightarrow axb<ayb.
\]
The group is left-orderable if it admits a left-invariant strict total order, and bi-orderable if it admits a strict total order invariant on both sides.

\textbf{Problem (A. I. Kokorin, 1965; Kourovka Notebook 1.51).} The supplied question is: \emph{What conditions ensure a matrix group over a field (of complex numbers) to be orderable?}

In the notation above, the requested task is to identify conditions on \(G\leq\operatorname{GL}_n(K)\) ensuring the existence of the intended invariant total order. The supplied record does not specify whether ``orderable'' means left-orderable or bi-orderable. It also leaves unclear whether \(K\) ranges over arbitrary fields, over subfields of \(\C\), or is specifically \(\C\). These alternatives are unresolved qualifications of the statement. The question does not specify whether the desired conditions are merely sufficient or constitute a necessary-and-sufficient characterization.
\subsection{Short English statement}
Kokorin asks for conditions ensuring that a matrix group admits an invariant total order. The supplied formulation leaves both the intended invariance of the order and the permitted fields unclear; it also does not specify whether a complete characterization is required.
\subsection{Sources}
\begin{itemize}
\item Local UnsolvedMath registry, permanent ID \(700046\), source problem code \texttt{catalogue entry}, title \emph{Kourovka Notebook Problem 1.51}. Supplied upstream revision: \texttt{d371ee917ac9647fbb06cd6fc529deec6f3dda82}.
\item Supplied source attribution: \emph{The Kourovka Notebook (Unsolved Problems in Group Theory)}, original problem 1.51, first issue (1965); proposed by A. I. Kokorin.
\item The supplied background identifies the twenty-first issue edition (2026), \texttt{arXiv:1401.0300v48}. The collection description credits E. I. Khukhro and V. D. Mazurov as editors. This attribution and edition information have not been independently verified here.
\item The local record labels the problem \emph{open}. This is a supplied registry status, not an independently verified assessment of the literature. No separate literature sources or literature-verification date are supplied.
\end{itemize}
\end{document}
